%% ---------------------------------------------------------------------------
%% Prescribed finite-window reaction-time modality by conserved-budget
%% support design.
%%
%% Single self-contained submission draft for
%% Journal of Physics A: Mathematical and Theoretical (IOP Publishing).
%%
%% Converted from the two-file APS RevTeX package
%%   encounter_multimodal_prr_submission.tex            (main text)
%%   + exact_m_theorem_spine.tex                        (\input)
%%   encounter_multimodal_prr_submission_supplement.tex (Supplemental Material)
%%   + exact_m_theorem_full_proof.tex                   (\input)
%% into ONE document: the former separate Supplemental Material file is now
%% appendices A-C of this article.  Accordingly there is no separate SM, the
%% SM placeholder bibliography entry has been dropped from references.bib,
%% and no citation to it remains anywhere in the text.
%%
%% Class: iopjournal.cls, IOP Publishing's current official LaTeX class
%% (2025 template, https://publishingsupport.iopscience.iop.org/questions/latex-template/).
%% It supersedes the older iopart.cls, which IOP no longer distributes.
%% Bibliography style: iopart-num (IOP numeric).
%%
%% Build:  latexmk -pdf encounter_multimodal_jpa.tex
%% ---------------------------------------------------------------------------

\documentclass{iopjournal}

% Options available in this class:
%   [anonymous]  suppress author names, affiliations and acknowledgments
%                for double-anonymous peer review.

\usepackage{cmap}
\usepackage[T1]{fontenc}
\usepackage{lmodern}
\usepackage{amsmath}
\usepackage{amssymb}
\usepackage{mathtools}
\usepackage{amsthm}
\usepackage{bm}
\usepackage{microtype}

% hyperref is loaded by iopjournal.cls; only adjust its settings here.
\hypersetup{
  pdftitle={Prescribed finite-window reaction-time modality by conserved-budget support design},
  pdfauthor={Xiaoxiao Zhouyi},
  pdfsubject={Exact finite-modality theorem for Doi encounter reactions},
  pdfkeywords={encounter reaction, reaction-time density, multimodality, conserved reactivity, Doi model}
}

%% ---- shared macros (union of the four source files) -----------------------
\newcommand{\dd}{\mathrm d}
\newcommand{\e}{\mathrm e}
\newcommand{\Lfree}{\mathcal L}
\newcommand{\Qbox}{\mathcal Q_d^{\mathrm{box}}}
\newcommand{\Qinf}{\mathcal Q_d^{\infty}}

%% ---- theorem environments, defined once for the whole document ------------
%% One shared counter; plain numbering in the main text (Theorem 1), and
%% section-qualified numbering inside the appendices (Lemma C.1, ...).
\theoremstyle{plain}
\newtheorem{theorem}{Theorem}
\newtheorem{lemma}[theorem]{Lemma}
\newtheorem{proposition}[theorem]{Proposition}
\newtheorem{corollary}[theorem]{Corollary}
\theoremstyle{definition}
\newtheorem{remark}[theorem]{Remark}

%% ---- section-number prefix, so that appendix headings read "Appendix A" ---
%% \secprefix is empty in the main text and set to "Appendix~" after \appendix.
%% Only \section-level headings are prefixed, so subsections stay "C.1" and
%% equation numbers stay "(C.1)".
\makeatletter
\newcommand{\secprefix}{}
\renewcommand{\@seccntformat}[1]{%
  \expandafter\ifx\csname #1\endcsname\section\secprefix\fi
  \csname the#1\endcsname\quad}
\makeatother

\date{\today}

\begin{document}

\articletype{Paper}

\title{Prescribed finite-window reaction-time modality by conserved-budget
support design}

% ORCID: https://orcid.org/0009-0003-7118-9621
\author{Xiaoxiao Zhouyi\orcid{0009-0003-7118-9621}}

\affil{School of Engineering Mathematics and Technology, University of Bristol,
Bristol BS8 1TW, United Kingdom}

\email{xiaoxiao.zhouyi@bristol.ac.uk}

\keywords{encounter reaction, reaction-time density, multimodality,
conserved reactivity, Doi model}

\begin{abstract}
Can a fixed amount of spatial reactivity be arranged to prescribe the
finite-window topology of an encounter-reaction-time density?  We answer this
existence question for a Doi encounter process.  For every fixed finite
physical dimension \(d\ge2\), prescribed fixed finite integer \(m\), and
compact family of interior catalyst weights, an \(m\)-dependent design of
ordered narrow catalyst slabs along a monotone exposure trajectory produces
exactly \(m\) nondegenerate maxima and \(m-1\) nondegenerate minima on a
declared compact positive-time window.  A global zero bound for the separated
Gaussian exposure clocks, a whole-window complement certificate, and
compact-positive-time mixed-jet convergence transfer the complete stationary
signature to the positive-budget killed process.  The construction first
fixes a sufficiently small noise-and-width scale and only then takes a
sufficiently small installed reactivity budget.  It is deliberately
nonuniform: the geometry and admissible scales may depend on \(d\) and \(m\),
and no useful budget threshold, event-mass floor, or stationary-point count
outside the declared window is obtained.  The theorem therefore establishes
prescribed finite modality under a conserved spatial resource without
asserting a particular finite-parameter two-dimensional realization.
\end{abstract}

\vspace{3mm}
\noindent{\fontsize{8}{10}\selectfont\raggedright
{\bfseries Draft compiled:} \today}\vspace{2mm} % placeholder date, replace at submission

\section{Introduction}
\label{sec:introduction}

First-passage and reaction-time densities may develop several peaks when
distinct trajectory families reach a reactive region on separated time
scales \cite{godecMetzler2017twochannel,
grebenkovMetzlerOshanin2018defocusing}.  Continuum trap placement can alter
capture-time topology \cite{lindsaySpoonmoreTzou2016narrow}, and
two-particle encounter laws in physical dimensions two and three show related
multirate mechanisms \cite{leVotEtAl2022encounter,giuggioli2020confined,
giuggioliEtAl2024multitarget}.  Spatially heterogeneous absorption and
fixed-total-reactivity optimization are also well established
\cite{nguyenGrebenkov2010porous,grebenkov2019heterogeneous,
nicolaouMulder2023flux}.  Here we ask a narrower existence question: can the
finite-window timing topology of a two-particle reaction be prescribed by
support design while the installed amount of static reactivity remains fixed?

The fact that a univariate mixture of \(m\) Gaussian components has at most
\(m\) modes is classical \cite{carreiraPerpinanWilliams2003modes,
amendolaEngstromHaase2020modes}; that upper bound is not our contribution.
The issue for an encounter process is to construct every stationary point
with a strict margin, exclude additional stationary points throughout a
declared time window after the contact factor is included, and carry the
complete signature from free exposure to a killed Doi law.

We prove that this can be done for every prescribed fixed finite \(m\) in
every fixed finite dimension \(d\ge2\).  The supports, contact geometry,
transport, initial law, and total installed centre-space reactivity are fixed
within each constructed family.  All interior allocations in that family
share the same alternating signature.  Thus the result concerns
support-geometry design under a conserved budget; it is not a statement that
one fixed support family switches among different mode counts when its
weights are reallocated.

The Doi killing representation, the associated reactivity-times-occupancy
identity, and Gaussian-mixture mode bounds are prior ingredients
\cite{doi1976stochastic,isaacsonMauroNewby2016uniform,
pruestelMeierSchellersheim2014area,bressloff2022feynmanKac}.  The new content
is their margin-bearing combination: a separated-clock construction with a
global zero count, a uniform whole-window sign certificate, stability under
the slowly varying contact probability, and a positive-budget transfer of
the full finite stationary signature.  The result concerns a transient
density under static spatial allocation and is distinct from universal
mean-first-passage response to rare activation
\cite{keidarReuveni2026linear}.

This article is self-contained.  Sections~\ref{sec:model}--\ref{sec:transfer}
fix the model and the positive-time transfer estimate,
section~\ref{sec:exact-m-spine} states the modality theorem and explains its
mechanism, and sections~\ref{sec:finite-scope}--\ref{sec:discussion} set out
the scope.  The analytic realizations of the generator, the mixed-jet bound
with its explicit constant, and the complete proof of the theorem are given
in appendices~\ref{app:doi-model}--\ref{app:full-proof}.

\section{Conserved-reactivity encounter model}
\label{sec:model}

Let \(q(x,t)\) be the unreacted configuration density of two mobile
particles, with conservative generator \(\mathcal L\).  Write \(r\) for the
relative coordinate and \(c\) for an encounter-centre coordinate.  For a
fixed contact profile \(\chi_a(r)\) and normalized nonnegative centre-space
profiles \(\phi_j(c)\), define
\begin{align}
 \kappa_{\bm w}(c)&=B\sum_{j=1}^{J}w_j\phi_j(c),
 &w_j&\ge0,\nonumber\\
 \sum_jw_j&=1,
 &\int\phi_j(c)\,\dd c&=1,\nonumber\\
 K_{B,\bm w}(x)&=\chi_a(r)\kappa_{\bm w}(c).
 \label{eq:budget-field}
\end{align}
Every admissible allocation has the same installed centre-space resource,
\begin{equation}
 \int\kappa_{\bm w}(c)\,\dd c=B.
 \label{eq:conserved-budget}
\end{equation}
This is not an integral over the complete two-particle configuration space.
The killed evolution, reaction-time density, and survival probability are
\begin{align}
 \partial_tq_{B,\bm w}&=(\mathcal L-K_{B,\bm w})q_{B,\bm w},
 &q_{B,\bm w}(0)&=q_0,\nonumber\\
 f_{B,\bm w}(t)&=\langle K_{B,\bm w},q_{B,\bm w}(t)\rangle,
 &-S'_{B,\bm w}(t)&=f_{B,\bm w}(t).
 \label{eq:killed-law}
\end{align}
Equation~\eqref{eq:killed-law} is the usual mass balance for a Doi
volume-reaction process, not a novelty claim.

For longitudinal catalyst slabs, translation invariance of the common
transverse coordinate gives an exact quotient
\begin{equation}
 \Omega_\infty
 =\mathbb R_z\times\mathbb R_{r_\parallel}\times
   \mathbb T_W^{d-1}.
 \label{eq:quotient}
\end{equation}
In the construction below the midpoint \(Z_t\) moves monotonically in mean,
the relative process remains a positive distance inside the contact region
along its deterministic trajectory on the chosen window, and each normalized
slab is centred at the midpoint mean at a target time.  The
\(W^{-(d-1)}\) factor in the killing field accounts for the omitted common
transverse directions and preserves Eq.~\eqref{eq:conserved-budget}.
Appendix~\ref{app:doi-model} records the bounded reflecting/periodic and
unbounded natural-decay realizations of this quotient together with the
analyticity and positivity properties of the killed semigroup that the
transfer estimate uses.

\section{Positive-time transfer}
\label{sec:transfer}

Write \(K_{B,\bm w}=BV_{\bm w}\) and introduce interior affine control
coordinates \(\bm w(\bm\theta)=\bar{\bm w}+P\bm\theta\).  The normalized
reaction density and the free-exposure law are
\begin{equation}
 F_B(t,\bm\theta)=B^{-1}f_{B,\bm w(\bm\theta)}(t),
 \qquad
 G(t,\bm\theta)=
 \langle V_{\bm w(\bm\theta)},\e^{t\mathcal L}q_0\rangle.
 \label{eq:normalized-exposure}
\end{equation}
On either a bounded reflecting/periodic quotient or the reversible
natural-decay realization of Eq.~\eqref{eq:quotient}, assume \(q_0\) belongs
to the corresponding state space and \(V_{\bm w}\) is bounded on a complex
control neighbourhood of a compact interior control set.  A
Dyson expansion followed by time and control Cauchy estimates gives, for
every fixed time order \(r\) and finite control multi-index \(\alpha\),
\begin{equation}
 \sup_{(t,\bm\theta)\in[\tau,T]\times\Theta}
 \left|
 \partial_t^r\partial_{\bm\theta}^{\alpha}(F_B-G)
 \right|
 \le C_{r,\alpha,\tau,T,\Theta}B.
 \label{eq:mixed-jet}
\end{equation}
Appendix~\ref{app:mixed-jet} gives the explicit constant and the proof.  In
particular, \(C^2\) convergence on a compact positive-time window preserves
disjoint stationary-point neighbourhoods, curvature signs, endpoint slopes,
and a nonzero derivative margin on their complement.  The estimate is local
in positive time and provides neither a useful numerical value of \(B\) nor
control on the scale \(t=O(B^{-1})\).

\section{Prescribed finite modality from separated exposure clocks}
\label{sec:exact-m-spine}

The conserved-budget support construction admits encounter processes with any
prescribed \emph{fixed finite} number of reaction-time modes.  We state this
existence result in an exact quotient of a two-particle Doi process; the
construction depends on the prescribed dimension and mode count and does not
assert mode-count switching by weight redistribution on one fixed support.

Fix finite $d\ge2$, finite $m\ge1$, a transverse period $W>0$, a contact
radius $0<a<W/2$, positive scales $\ell_0,\rho>0$, and a compact interval
$I=[\tau,T]\subset(0,\infty)$.
On $\mathbb R\times\mathbb R\times\mathbb T_W^{d-1}$, consider the
independent midpoint and relative processes

\begin{equation}
 \begin{aligned}
 dZ_t&=-\gamma(Z_t-\bar z)\,dt
       +\varepsilon\sqrt{D_0}\,dW_t,\\
 dR_{\parallel,t}&=-\gamma R_{\parallel,t}\,dt
       +2\varepsilon\sqrt{D_0}\,dW_{\parallel,t},\\
 dR_{\perp,t}&=2\varepsilon\sqrt{D_0}\,dW_{\perp,t}\pmod W.
 \end{aligned}
 \label{eq:exact-m-process}
\end{equation}

The mutually independent initial laws satisfy

\begin{equation}
 \begin{aligned}
 Z_0&\sim N\!\left(z_0,\varepsilon^2D_0/(2\gamma)\right),\\
 R_{\parallel,0}&\sim N(r_{\parallel,0},\varepsilon^2u_0^2),\\
 R_{\perp,0}&\sim
 \operatorname{WN}(r_{\perp,0},\varepsilon^2\Sigma_{\perp,0}),
 \end{aligned}
 \qquad
 \begin{gathered}
 D_0,\gamma,u_0>0,\quad z_0\ne\bar z,\\
 u_0^2<4D_0/\gamma,\quad \Sigma_{\perp,0}\succ0.
 \end{gathered}
 \label{eq:exact-m-initial}
\end{equation}

Here $\operatorname{WN}$ is the wrapped Gaussian on
$\mathbb T_W^{d-1}$, $\Sigma_{\perp,0}$ is a positive-definite
$(d-1)\times(d-1)$ covariance matrix, and $|r|_{\rm mi}$ is the Euclidean norm
of the longitudinal relative coordinate together with the shortest transverse
representatives modulo $W$.

All driving Brownian motions are also independent.  Thus
$\mu(t)=\bar z+(z_0-\bar z)e^{-\gamma t}$ is strictly monotone and
$\operatorname{Var}Z_t=\varepsilon^2D_0/(2\gamma)$.  Under the minimum-image
convention, assume the deterministic relative trajectory
$r_*(t)=(r_{\parallel,0}e^{-\gamma t},r_{\perp,0})$ obeys
$\sup_{t\in I}|r_*(t)|_{\rm mi}\le a-\eta$ for some $\eta>0$.

Choose target times

\begin{equation}
 \begin{aligned}
  &\tau<t_1<\cdots<t_m<T,\qquad c_j=x(t_j),\\
  &x(t)=\frac{\operatorname{sgn}(\mu')\mu(t)}{\ell_0},
 \end{aligned}
 \label{eq:exact-m-centres}
\end{equation}

so that $x'(t)$ is bounded away from zero.  At the corresponding physical
centres, place normalized Gaussian catalyst slabs of width
$\varepsilon\rho$.  Their nonnegative weights belong to a compact subset of
the simplex,

\begin{equation}
 \mathcal W_{w_*}
 =\left\{\bm w:\
   \substack{\sum_{j=1}^{m}w_j=1,\\ w_j\ge w_*>0}\right\},
 \qquad 0<w_*\le\frac1m,
 \label{eq:exact-m-simplex}
\end{equation}

and define

\begin{equation}
 \begin{aligned}
 \phi_{j,\varepsilon}(z)
 &=\frac{\exp[-(z-\mu(t_j))^2/(2\varepsilon^2\rho^2)]}
        {\sqrt{2\pi}\,\varepsilon\rho},\\
 K_{B,\bm w,\varepsilon}(Z,R)
 &=\bm1_{\{|R|_{\rm mi}<a\}}\frac{B}{W^{d-1}}
   \sum_{j=1}^{m}w_j\phi_{j,\varepsilon}(Z).
 \end{aligned}
 \label{eq:exact-m-killing}
\end{equation}

The slab normalization makes $B$ the same installed centre-space budget for
every allocation.  Set $K_{B,\bm w,\varepsilon}=B V_{\bm w,\varepsilon}$,
and let $T_0(t)$ and $q_0$ denote the unkilled semigroup and initial law.  Then

\begin{equation}
 \begin{aligned}
 G_{\varepsilon,\bm w}(t)
 &=B^{-1}\mathbb E_0[K_{B,\bm w,\varepsilon}(Z_t,R_t)]
 =\langle V_{\bm w,\varepsilon},T_0(t)q_0\rangle,\\
 f_{B,\varepsilon,\bm w}(t)
 &=\mathbb E_0\!\left[K_{B,\bm w,\varepsilon}(Z_t,R_t)
  e^{-\int_0^tK_{B,\bm w,\varepsilon}(Z_s,R_s)\,ds}\right].
 \end{aligned}
 \label{eq:exact-m-observables}
\end{equation}

The expectation $\mathbb E_0$ is with respect to the unkilled process started
from Eq.~\eqref{eq:exact-m-initial}.  Independence gives the exact
factorization

\begin{align}
 G_{\varepsilon,\bm w}(t)
 &=A_\varepsilon(t)H_{\sigma,\bm w}(x(t)),\nonumber\\
 A_\varepsilon(t)
 &=\frac{c_{d,\varepsilon}(t)}
 {W^{d-1}\sqrt{2\pi}\,\varepsilon
  \sqrt{D_0/(2\gamma)+\rho^2}},
 \label{eq:exact-m-factorization}
\end{align}

where $c_{d,\varepsilon}(t)=
\mathbb P_0(|R_t|_{\rm mi}<a)$ and $H_{\sigma,\bm w}$ is defined below.

\begin{theorem}[Exact prescribed finite modality]
\label{thm:main-modality}
Fix finite $d\ge2$, finite $m\ge1$, the data above, and the compact window
$I$.  There is $\varepsilon_0>0$ such that, for every
$0<\varepsilon<\varepsilon_0$ and every
$\bm w\in\mathcal W_{w_*}$, the unit-budget free-exposure clock has exactly
$m$ nondegenerate maxima and $m-1$ nondegenerate minima on $I$, ordered
alternately, and no stationary endpoint.  For each fixed such
$\varepsilon$, there is $B_0(\varepsilon)>0$, uniform over
$\mathcal W_{w_*}$, such that the budget-rescaled Doi reaction-time density
$f_{B,\varepsilon,\bm w}/B$ has the same complete stationary signature for
every $0<B<B_0(\varepsilon)$.  Multiplication by $B$ leaves the stationary
points unchanged, so the conclusion also holds for the physical density
$f_{B,\varepsilon,\bm w}$.
\end{theorem}

\noindent
Theorem~\ref{thm:main-modality} is proved in
appendix~\ref{app:full-proof}, where it is restated with its complete
hypothesis list as theorem~\ref{thm:exactmfull-doi}.  The remainder of this
section explains the mechanism.

The mechanism can be seen without invoking a normal form.  Gaussian
convolution reduces the longitudinal part of the unit-budget exposure to

\begin{equation}
 \begin{aligned}
  H_{\sigma,\bm w}(x)
  &=\sum_{j=1}^{m}w_j
    \exp\!\left[-\frac{(x-c_j)^2}{2\sigma^2}\right],\\
  \sigma&=\frac{\varepsilon}{\ell_0}
  \sqrt{\frac{D_0}{2\gamma}+\rho^2}.
 \end{aligned}
 \label{eq:exact-m-mixture}
\end{equation}

Writing $q_j(x)=w_j\exp[-(x-c_j)^2/(2\sigma^2)]$,
$\pi_j(x)=q_j(x)/H_{\sigma,\bm w}(x)$, and
$L=\partial_x\log H$, let $\bar c(x)$ and
$\operatorname{Var}_{\pi(x)}(c)$ denote the mean and variance of the centres
under the finite probability vector $\bm\pi(x)$.  Then

\begin{equation}
 L(x)=\frac{\bar c(x)-x}{\sigma^2},
 \qquad
 L'(x)=\frac{\operatorname{Var}_{\pi(x)}(c)}{\sigma^4}
       -\frac{1}{\sigma^2}.
 \label{eq:exact-m-log-slope}
\end{equation}

The univariate upper bound of $m$ modes for an $m$-component Gaussian mixture
is classical \cite{carreiraPerpinanWilliams2003modes,
amendolaEngstromHaase2020modes}.  Here a direct affine--exponential proof is
retained because the Doi transfer needs all critical zeros counted with
multiplicity together with uniform root-tube, complement, curvature, and
endpoint margins.

For the adjacent pair $c_j<c_{j+1}$, the weighted crossover is

\begin{equation}
 s_j=\frac{c_j+c_{j+1}}{2}
 +\frac{\sigma^2}{c_{j+1}-c_j}
  \log\!\frac{w_j}{w_{j+1}}.
 \label{eq:exact-m-crossover}
\end{equation}

The adjacent ratio obeys
$q_{j+1}/q_j=1/9$ at
$s_j-\sigma^2\log(9)/(c_{j+1}-c_j)$ and $q_{j+1}/q_j=9$ at the corresponding
plus edge.  Thus both adjacent components remain uniformly present throughout
the crossover, while all nonadjacent components are exponentially small.
Equation~\eqref{eq:exact-m-log-slope} then yields one simple minimum at

\begin{equation}
 r_j=s_j+O(\sigma^4),
 \qquad
 L'(r_j)=\frac{(c_{j+1}-c_j)^2}{4\sigma^4}
          +O(\sigma^{-2}).
 \label{eq:exact-m-valley}
\end{equation}

Each centre similarly contributes one simple maximum at
$p_j=c_j+O(\exp[-q/\sigma^2])$, with
$L'(p_j)=-\sigma^{-2}+o(\sigma^{-2})$.  These constructions give
$2m-1$ simple stationary points.  They are exhaustive: after multiplication
by a positive Gaussian factor, $H'$ is an exponential polynomial

\begin{equation}
 \sum_{j=1}^{m}(a_j+b_jx)e^{\lambda_jx},
 \qquad \lambda_1<\cdots<\lambda_m,
 \label{eq:exact-m-exp-polynomial}
\end{equation}

which has at most $2m-1$ real zeros counted with multiplicity.  The extreme
affine--exponential terms dominate in the two tails, so all zeros lie in a
compact interval and real analyticity makes their number finite.  A
generalized Rolle induction then proves the bound by differentiating twice to
remove the lowest-exponent affine term.

The remaining relative-coordinate contact probability is a positive slow
factor.  Uniform bounds on its first two logarithmic time derivatives, peak
boxes of width $O(\sigma^2)$, valley boxes of width $O(\sigma^4)$, and a
posterior-sector sign certificate exclude stationary points on the full box
complement.  Consequently the slow factor moves maxima by $O(\sigma^2)$ and
minima by $O(\sigma^4)$ without changing their number or type.  Finally, on
every fixed positive $\varepsilon$, compact-positive-time mixed-jet
convergence

\begin{equation}
 \sup_{\bm w\in\mathcal W_{w_*}}
 \left\|\frac{f_{B,\varepsilon,\bm w}}{B}
       -G_{\varepsilon,\bm w}\right\|_{C^2(I)}
 \longrightarrow0
 \qquad (B\downarrow0)
 \label{eq:exact-m-weak-budget}
\end{equation}

preserves the disjoint root tubes, signed curvature margins, complement
margin, and endpoint slopes.

The quantifier order is essential: the geometry and admissible scales may
depend on fixed $(d,m)$; one first fixes a sufficiently small
$\varepsilon>0$ and only then chooses $B<B_0(\varepsilon)$.  The theorem does
not provide a useful uniform lower bound for $B_0$, a uniform-in-$d$ or
uniform-in-$m$ construction, an event-mass floor, or a stationary-point count
outside $I$; it also uses a whole-window contact-interior condition whose
small-noise limit is asymptotically saturated contact.  The proof permits
$B_0(\varepsilon)$ to shrink rapidly as $\varepsilon$ decreases---potentially
exponentially in $1/\varepsilon^2$---and no rate is claimed; this is why no
common useful budget is asserted across the family.  Nontrivial contact,
observability at a common positive budget, and finite-parameter robustness
remain separate numerical validation gates.

\section{Finite-parameter scope}
\label{sec:finite-scope}

The theorem is complete without a discretized example.  Its proof establishes
an exact existence statement for the continuum Doi process: it uses no
finite-volume approximation, and no discretized convergence statement is
needed for its conclusion.  What the sequential small-parameter construction
does not do is identify a practically resolved physical parameter set.

Accordingly, this article includes no numerical evidence for a particular
finite-parameter realization, no off-lattice validation, no phase diagram,
and no finite-parameter dimensional comparison.  Its scientific conclusions
rest only on the exact theorem.

We also make no rigorous numerical claim that a particular finite-volume
scheme converges, together with its stationary roots, from bounded grids to
the unbounded continuum process.  Such a stronger numerical statement would
require quantitative consistency, spatial and box-truncation error bounds,
and componentwise margin-preserving transfer of each stationary component.
Those additional ingredients are unnecessary for the theorem proved here and
do not limit its scope.

%% =========================================================================
%% OPTIONAL: Numerical illustration (pending Isambard campaign)
%% -------------------------------------------------------------------------
%% Uncomment and populate this block when the Isambard numerical campaign is
%% complete.  It is deliberately placed after "Finite-parameter scope" and
%% before "Discussion", so that the caveats in section~\ref{sec:finite-scope}
%% can be relaxed in the same edit that adds the supporting evidence.
%% Remember to (i) revise section~\ref{sec:finite-scope}, (ii) add the figure
%% files to this directory, and (iii) add a real data-availability statement
%% in place of the purely-mathematical one below.
%%
%\section{Numerical illustration}
%\label{sec:numerics}
%
%TODO: describe the discretization, the parameter set, and the measured
%stationary signature; state explicitly which theorem hypotheses are
%satisfied numerically and which are only approximated.
%
%\begin{figure}
%  \centering
%  \includegraphics[width=0.8\textwidth]{FIGURE_FILENAME}
%  \caption{\label{fig:numerics}TODO caption.}
%\end{figure}
%% =========================================================================

\section{Discussion}
\label{sec:discussion}

The construction turns catalyst locations into separated exposure clocks.
Narrow Gaussian convolution creates one peak near each target time; an
affine-exponential zero bound makes the count exhaustive; posterior-sector
estimates control the complement after multiplication by the contact
probability; and the positive-time estimate
Eq.~\eqref{eq:mixed-jet} transfers the strict topology to the Doi law.
Together these steps establish arbitrary prescribed fixed finite modality
under a conserved centre-space reactivity budget.

The restrictions are substantive.  Geometry and admissible scales depend on
\(d\) and \(m\).  The deterministic relative trajectory stays strictly
inside the contact set on the design window, and its free contact probability
approaches one there as the noise scale vanishes.  The limits are sequential,
observability at a common useful budget is not guaranteed, and nothing is
proved about topology outside the declared window.  These limits distinguish
the exact existence theorem from a finite-parameter design rule while leaving
its finite-window conclusion intact.

%TODO(author): decide whether an Acknowledgments and/or Funding section is
%needed before submission.  If so, uncomment and complete the two commands
%below.  JPA requires a funding statement only where funding was received;
%if there was none, leave both commented out.
%\ack{TODO}
%\funding{TODO}

\data{This is a purely mathematical work.  No new data were created or
analysed in this study, and therefore no data-availability statement beyond
this one is applicable.}

%% ===========================================================================
%% APPENDICES
%% These absorb, in full, the former Supplemental Material of the APS version:
%%   appendix A  = SM section S1  (Doi model and analytic realizations)
%%   appendix B  = SM section S2  (uniform mixed-jet theorem and proof)
%%   appendix C  = exact_m_theorem_full_proof.tex (complete proof)
%% The SM's own "Scope of the analytical result" section has been merged into
%% section~\ref{sec:finite-scope} of the main text and is not repeated here.
%% ===========================================================================

\appendix

%% From here on, equations and theorem-like environments are numbered by
%% appendix letter: (A.1), (B.1), Lemma C.1, ...
\numberwithin{equation}{section}
\numberwithin{theorem}{section}
\makeatletter\renewcommand{\secprefix}{Appendix~}\makeatother

\section{Doi model and analytic realizations}
\label{app:doi-model}

For a fixed finite physical dimension \(d\ge2\), the longitudinal-slab
symmetry reduction gives
\begin{equation}
 \Qbox=I_z\times I_\parallel\times\mathbb T_W^{d-1},
 \qquad
 \Qinf=\mathbb R\times\mathbb R\times\mathbb T_W^{d-1}.
 \label{eq:supp-domains}
\end{equation}
With \(x=(z,r_\parallel,r_\perp)\), equal particle diffusivities give the free
forward operator
\begin{equation}
 \Lfree q=-\nabla\!\cdot(bq)+\nabla\!\cdot(\bm D\nabla q),
 \qquad
 \bm D=\operatorname{diag}(D/2,2D,\ldots,2D),
 \label{eq:supp-free-generator}
\end{equation}
where
\begin{equation}
 b(x)=\bigl(-\gamma(z-\bar z),-\gamma r_\parallel,0,\ldots,0\bigr).
 \label{eq:supp-drift}
\end{equation}
The bounded quotient has zero forward flux on its nonperiodic faces and
periodic transverse coordinates.  Natural decay on the unbounded quotient is
specified by the reversible density
\begin{equation}
 \pi(x)=Z_\pi^{-1}
 \exp\!\left[-\frac{\gamma(z-\bar z)^2}{D}
              -\frac{\gamma r_\parallel^2}{4D}\right],
 \label{eq:supp-pi}
\end{equation}
including the normalized uniform torus factor.

Let \(0<a<W/2\), let \(\chi_a(r)\) be the embedded minimum-image contact-ball
indicator, and let normalized nonnegative longitudinal profiles obey
\(\phi_j\in L^\infty\) and \(\int\phi_j(z)\,\dd z=1\).  Put
\begin{equation}
 V_j(z,r)=W^{-(d-1)}\chi_a(r)\phi_j(z),\qquad
 V_w=\sum_{j=1}^Jw_jV_j,\qquad K_{B,w}=BV_w,
 \label{eq:supp-profiles}
\end{equation}
where \(w_j\ge0\) and \(\sum_jw_j=1\).  The installed centre-space catalyst
amount is \(B\) for every \(w\).  The killed state and normalized reaction
density are
\begin{equation}
 q_{B,w}(t)=\e^{t(\Lfree-BM_{V_w})}q_0,\qquad
 F_B(t,w)=\langle V_w,q_{B,w}(t)\rangle .
 \label{eq:supp-state}
\end{equation}

\begin{proposition}[Analytic killed semigroups]
\label{prop:supp-analytic}
On \(\Qbox\), the reflecting/periodic realization of \(\Lfree\) on
\(L^2(\Qbox)\) generates a positive analytic semigroup.  On \(\Qinf\), the
natural-decay realization on \(X_\pi=L^2(\pi^{-1}\dd x)\) does the same.
For real \(B\ge0\) and simplex control \(w\), bounded multiplication by
\(-BV_w\) preserves the generator domain and yields a positive,
mass-decreasing analytic semigroup.  At positive time the state is entire in
every finite affine complexification of the control amplitudes.  On
\(\Qinf\), the observable pairing is continuous when
\(V_w\in L^2(\pi\dd x)\cap L^\infty\).
\end{proposition}

\begin{proof}
Set \(q=\pi u\).  The identity
\begin{equation}
 \Lfree(\pi u)=\pi\mathcal Gu,\qquad
 \mathcal G=\pi^{-1}\nabla\!\cdot(\pi\bm D\nabla)
 \label{eq:supp-similarity}
\end{equation}
follows from \(b=\bm D\nabla\log\pi\).  On \(L^2(\pi\dd x)\), the closed
weighted form has
\begin{equation}
 \langle u,\mathcal Gv\rangle_{L^2(\pi)}
 =-\int(\nabla u)^{\mathsf T}\bm D\nabla v\,\pi\dd x.
 \label{eq:supp-form}
\end{equation}
It therefore realizes \(\mathcal G\) as a nonpositive self-adjoint generator.
Multiplication by \(\pi\) is a bounded isomorphism on the bounded quotient and
a unitary map into \(X_\pi\) on the unbounded quotient.  Bounded multiplication
by \(-BV_w\) preserves analyticity and the domain.  The norm-convergent
Dyson--Phillips series gives entire dependence on the finite control
amplitudes.  Positivity and mass decrease for real nonnegative data follow
from the product formula, equivalently from Feynman--Kac
\cite{bressloff2022feynmanKac,grebenkov2020multiple}.
\end{proof}

\section{Uniform compact-positive-time mixed derivatives}
\label{app:mixed-jet}

Choose an interior affine control chart \(w(\theta)=w_*+P\theta\) and a
compact real control set \(\Theta\).  Fix \(\delta>0\) such that each closed
coordinate polydisc of radius \(\delta\) about a point of \(\Theta\) lies in
the complex tube
\begin{equation}
 \Theta_\delta^+
 =\{\zeta\in\mathbb C^{J-1}:
   \operatorname{dist}_\infty(\zeta,\Theta)<2\delta\}.
 \label{eq:supp-tube}
\end{equation}
Define
\begin{equation}
 G(t,\theta)=\langle V_{w(\theta)},\e^{t\Lfree}q_0\rangle,\qquad
 v_{\infty,\delta}=
 \sup_{\zeta\in\Theta_\delta^+}\|V_{w(\zeta)}\|_\infty .
 \label{eq:supp-free-exposure}
\end{equation}
Let \(v_{*,\delta}\) be the supremum of
\(\|V_{w(\zeta)}\|_{L^2}\) on the bounded quotient and of
\(\|V_{w(\zeta)}\|_{L^2(\pi\dd x)}\) on the unbounded quotient.  Write
\(\kappa_\pi\) for the bounded-quotient similarity bound and set
\(\kappa_\pi=1\) on the unbounded quotient.

\begin{theorem}[Uniform weak-reaction mixed-jet bridge]
\label{thm:mixed-jet}
Fix \(0<\tau\le T<\infty\), \(0\le B\le B_{\max}\), an integer \(r\ge0\),
and a finite control multi-index \(\alpha\).  The normalized density
\(F_B(t,\theta)\) in Eq.~\eqref{eq:supp-state} extends continuously to
\(F_0=G\), and
\begin{align}
 &\sup_{(t,\theta)\in[\tau,T]\times\Theta}
 \left|\partial_t^r\partial_\theta^\alpha(F_B-G)\right|
 \nonumber\\
 &\quad\le
 r!\,\alpha!\left(\frac2\tau\right)^r
 \delta^{-|\alpha|}v_{*,\delta}\kappa_\pi
 \left[\exp\!\left(\frac32Bv_{\infty,\delta}T\right)-1\right]
 \|q_0\|_X .
 \label{eq:supp-mixed-bound}
\end{align}
In particular, the left side is at most \(BC_{r,\alpha}\) for a finite
constant independent of \(B\in[0,B_{\max}]\).  This is the explicit form of
Eq.~\eqref{eq:mixed-jet}.
\end{theorem}

\begin{proof}
For complex time \(z\) with \(\operatorname{Re}z>0\), the
norm-convergent Dyson expansion in the self-adjoint representation is
\begin{align}
 \e^{z(\mathcal G-BV_w)}
 ={}&\e^{z\mathcal G}
 +\sum_{n=1}^{\infty}(-Bz)^n
 \int_{\Delta_n}\e^{z(1-s_1)\mathcal G}V_w
 \left[\prod_{k=1}^{n-1}
 \e^{z(s_k-s_{k+1})\mathcal G}V_w\right]
 \e^{zs_n\mathcal G}\,\dd\bm s ,
 \label{eq:supp-dyson}
\end{align}
where
\(\Delta_n=\{1\ge s_1\ge\cdots\ge s_n\ge0\}\) has volume \(1/n!\).
Every free factor is contractive.  Consequently,
\begin{equation}
 \|\e^{z(\Lfree-BV_w)}-\e^{z\Lfree}\|_{X\to X}
 \le\kappa_\pi\{\e^{B|z|\|V_w\|_\infty}-1\}.
 \label{eq:supp-operator-bound}
\end{equation}
The complex-bilinear observable remainder thus satisfies
\begin{equation}
 \left|\left\langle V_{w(\zeta)},
 [\e^{z(\Lfree-BV_{w(\zeta)})}-\e^{z\Lfree}]q_0
 \right\rangle\right|
 \le v_{*,\delta}\kappa_\pi
 \{\e^{Bv_{\infty,\delta}|z|}-1\}\|q_0\|_X .
 \label{eq:supp-remainder}
\end{equation}
For \(t\in[\tau,T]\), the disk \(|z-t|\le\tau/2\) remains in the open
right half-plane and has \(|z|\le3T/2\).  Cauchy's estimate on this disk and
on each radius-\(\delta\) control polydisc proves
Eq.~\eqref{eq:supp-mixed-bound}.  Finally,
\(\e^{Bx}-1\le Bx\e^{B_{\max}x}\) gives the linear bound.
\end{proof}

The estimate transfers any fixed finite collection of strict
stationary-point neighbourhood, curvature, endpoint, and complement
inequalities.  It is asserted only on compact positive-time intervals and
does not supply a useful numerical budget threshold.

\section{Exact prescribed finite modality: complete proof}
\label{app:full-proof}

This appendix proves the exact finite-mode statement announced as
theorem~\ref{thm:main-modality} in section~\ref{sec:exact-m-spine}.  The
construction is pointwise in a fixed finite dimension and a fixed finite mode
count.  Its two small parameters are ordered: first the narrow-noise
parameter is fixed sufficiently small, and only then is the Doi budget taken
sufficiently small.  All stationary-point counts below are on the declared
compact positive-time window.

\subsection{Exact quotient, stationary midpoint variance, and whole-window
contact}
\label{subsec:exactmfull-physical}

Fix finite integers \(d\ge2\) and \(m\ge1\), a transverse period \(W>0\), a
contact radius \(0<a<W/2\), and a compact positive-time interval
\(I=[\tau,T]\subset(0,\infty)\).  On
\(\mathbb R\times\mathbb R\times\mathbb T_W^{d-1}\), let the longitudinal
midpoint and relative coordinates solve
\begin{align}
 \dd Z_t&=-\gamma(Z_t-\bar z)\,\dd t
       +\varepsilon\sqrt{D_0}\,\dd W_t,
 \label{eq:exactmfull-midpoint-sde}\\
 \dd R_{\parallel,t}&=-\gamma R_{\parallel,t}\,\dd t
       +2\varepsilon\sqrt{D_0}\,\dd W_{\parallel,t},
 \label{eq:exactmfull-relative-long-sde}\\
 \dd R_{\perp,t}&=2\varepsilon\sqrt{D_0}\,\dd W_{\perp,t}
       \pmod W,
 \label{eq:exactmfull-relative-perp-sde}
\end{align}
where \(D_0,\gamma>0\).  All Brownian drivers and all initial variables are
mutually independent, and
\begin{align}
 Z_0&\sim N\!\left(z_0,\varepsilon^2\frac{D_0}{2\gamma}\right),
 &z_0&\ne\bar z,
 \label{eq:exactmfull-midpoint-initial}\\
 R_{\parallel,0}&\sim
 N(r_{\parallel,0},\varepsilon^2u_0^2),
 &0<u_0^2&<\frac{4D_0}{\gamma},
 \label{eq:exactmfull-relative-long-initial}\\
 R_{\perp,0}&\sim
 \operatorname{WN}(r_{\perp,0},\varepsilon^2\Sigma_{\perp,0}),
 &\Sigma_{\perp,0}&\succ0.
 \label{eq:exactmfull-relative-perp-initial}
\end{align}
Here \(\operatorname{WN}\) denotes the wrapped Gaussian on the transverse
torus.  The two strict variance inequalities
\(D_0/(2\gamma)<D_0/\gamma\) and
\(u_0^2<4D_0/\gamma\) are precisely the unbounded Gaussian conditions that
place this initial law in the weighted state space used by
theorem~\ref{thm:mixed-jet}; the wrapped factor is square integrable against
the uniform torus measure.

The midpoint has mean and variance
\begin{equation}
 \mu(t)=\bar z+(z_0-\bar z)\e^{-\gamma t},
 \qquad
 \operatorname{Var}Z_t=\varepsilon^2\frac{D_0}{2\gamma}
 \quad(t\ge0).
 \label{eq:exactmfull-stationary-variance}
\end{equation}
Thus only the variance, not the mean, is stationary.  This constant variance
is what produces a common-variance Gaussian mixture below.

Let \(\varsigma=\operatorname{sgn}\mu'\in\{-1,1\}\), fix a physical
reference length \(\ell_0>0\), and define
\begin{equation}
 x(t)=\frac{\varsigma\mu(t)}{\ell_0},
 \qquad v_0:=\inf_{t\in I}x'(t)>0.
 \label{eq:exactmfull-increasing-coordinate}
\end{equation}
Choose target times and their dimensionless centres by
\begin{equation}
 \tau<t_1<\cdots<t_m<T,
 \qquad c_j=x(t_j).
 \label{eq:exactmfull-targets}
\end{equation}
Then \(c_1<\cdots<c_m\), and the physical centre of slab \(j\) is
\(\varsigma\ell_0c_j=\mu(t_j)\).  In particular, reversing the trajectory
orientation reverses the physical centres at the same time; no centre is
held fixed under the coordinate reversal.

Write
\begin{equation}
 r_*(t)=(r_{\parallel,0}\e^{-\gamma t},r_{\perp,0})
 \label{eq:exactmfull-relative-deterministic-path}
\end{equation}
and impose the whole-window, minimum-image contact margin
\begin{equation}
 \sup_{t\in I}|r_*(t)|_{\rm mi}\le a-\eta
 \quad\text{for some }\eta>0.
 \label{eq:exactmfull-whole-window-contact}
\end{equation}
The word ``whole-window'' is essential: the condition is required on every
point of \(I\), not merely near the target times.

\begin{lemma}[Whole-window differentiated contact tail]
\label{lem:exactmfull-contact-tail}
Let
\begin{equation}
 c_{d,\varepsilon}(t)
 =\mathbb P\{|R_t|_{\rm mi}<a\}.
 \label{eq:exactmfull-contact-probability}
\end{equation}
Under Eqs.~\eqref{eq:exactmfull-relative-long-sde}--
\eqref{eq:exactmfull-whole-window-contact}, for \(r=0,1,2\) there are finite
constants \(C_{r,d},N_{r,d},q_d>0\), independent of all sufficiently small
\(\varepsilon\), such that
\begin{equation}
 \left\|\partial_t^r(c_{d,\varepsilon}-1)\right\|_{L^\infty(I)}
 \le C_{r,d}\varepsilon^{-N_{r,d}}
       \exp(-q_d/\varepsilon^2).
 \label{eq:exactmfull-contact-tail-bound}
\end{equation}
Consequently, \(c_{d,\varepsilon}\ge1/2\) for sufficiently small
\(\varepsilon\), and the first two logarithmic derivatives of
\(c_{d,\varepsilon}\) are uniformly bounded on \(I\) and converge to zero
faster than any power of \(\varepsilon\).
\end{lemma}

\begin{proof}
In an unwrapped transverse chart the relative Gaussian has mean \(r_*(t)\)
and covariance \(\varepsilon^2\Sigma_R(t)\), with longitudinal and transverse
blocks
\begin{align}
 u_0^2\e^{-2\gamma t}
 +\frac{2D_0}{\gamma}(1-\e^{-2\gamma t}),
 \qquad
 \Sigma_{\perp,0}+4D_0t I_{d-1}.
 \label{eq:exactmfull-relative-covariance}
\end{align}
Because \(I\subset(0,\infty)\), \(u_0>0\), and
\(\Sigma_{\perp,0}\succ0\), this covariance and its first two time
derivatives are uniformly bounded on \(I\), and its eigenvalues lie between
two positive constants.  The reverse triangle inequality and
Eq.~\eqref{eq:exactmfull-whole-window-contact} give
\begin{equation}
 \inf_{t\in I}\inf_{|r|_{\rm mi}\ge a}
 d_{\rm cyl}(r,r_*(t))\ge\eta.
 \label{eq:exactmfull-contact-separation}
\end{equation}
The assumption \(a<W/2\) keeps the contact ball away from the transverse cut
locus.  Express the wrapped density as its Gaussian lattice-image sum.
Every time derivative of order at most two is the same Gaussian multiplied
by a polynomial in the scaled displacement and in \(\varepsilon^{-1}\),
with degree bounded independently of \(t\).  Equation
\eqref{eq:exactmfull-contact-separation}, the uniform covariance bounds, and
a standard Gaussian tail estimate therefore give, after summing the lattice
images,
\begin{equation}
 \sup_{t\in I}\int_{|r|_{\rm mi}\ge a}
 |\partial_t^r p_{\varepsilon,t}(r)|\,\dd r
 \le C_{r,d}\varepsilon^{-N_{r,d}}
       \e^{-q_d/\varepsilon^2},
 \qquad r=0,1,2.
 \label{eq:exactmfull-density-tail}
\end{equation}
Differentiation under the integral is justified by the same integrable
majorant and yields Eq.~\eqref{eq:exactmfull-contact-tail-bound}.  For small
\(\varepsilon\), its \(r=0\) case gives \(c_{d,\varepsilon}\ge1/2\), while
\begin{align}
 |\partial_t\log c_{d,\varepsilon}|
 &\le2|\partial_t c_{d,\varepsilon}|,\nonumber\\
 |\partial_t^2\log c_{d,\varepsilon}|
 &\le2|\partial_t^2c_{d,\varepsilon}|
      +4|\partial_tc_{d,\varepsilon}|^2.
 \label{eq:exactmfull-log-contact-bounds}
\end{align}
This proves the logarithmic-derivative statement.
\end{proof}

Fix \(\rho>0\), put
\begin{equation}
 S_*^2=\frac{D_0}{2\gamma}+\rho^2,
 \qquad \sigma=\frac{\varepsilon S_*}{\ell_0},
 \label{eq:exactmfull-common-width}
\end{equation}
and use the normalized physical slabs
\begin{equation}
 \phi_{j,\varepsilon}(z)
 =\frac{1}{\sqrt{2\pi}\,\varepsilon\rho}
  \exp\!\left[-\frac{(z-\mu(t_j))^2}
  {2\varepsilon^2\rho^2}\right].
 \label{eq:exactmfull-slabs}
\end{equation}
Let \(0<w_*\le1/m\) and
\begin{equation}
 \mathcal W_{w_*}
 =\left\{w=(w_1,\ldots,w_m):
   \sum_{j=1}^m w_j=1,\quad w_j\ge w_*>0\right\}.
 \label{eq:exactmfull-weight-set}
\end{equation}
For \(w\in\mathcal W_{w_*}\), define the Doi killing field
\begin{equation}
 K_{B,w,\varepsilon}(Z,R)
 =\mathbf{1}_{\{|R|_{\rm mi}<a\}}
  \frac{B}{W^{d-1}}\sum_{j=1}^m
  w_j\phi_{j,\varepsilon}(Z)
 =B V_{w,\varepsilon}(Z,R).
 \label{eq:exactmfull-killing}
\end{equation}
Every slab has unit full longitudinal integral, so \(B\) is the same
installed centre-space budget for every allocation.

\begin{lemma}[Exact common-variance free-exposure factorization]
\label{lem:exactmfull-factorization}
Let \(T_0(t)\) be the unkilled semigroup and \(q_0\) the initial law above.
The unit-budget free-exposure clock is exactly
\begin{align}
 G_{\varepsilon,w}(t)
 &:=\langle V_{w,\varepsilon},T_0(t)q_0\rangle
 =a_\varepsilon(t)H_{\sigma,w}(x(t)),
 \label{eq:exactmfull-factorization}\\
 a_\varepsilon(t)
 &:=\frac{c_{d,\varepsilon}(t)}
 {W^{d-1}\sqrt{2\pi}\,\varepsilon S_*},
 \label{eq:exactmfull-slow-factor}\\
 H_{\sigma,w}(x)
 &:=\sum_{j=1}^m w_j
 \exp\!\left[-\frac{(x-c_j)^2}{2\sigma^2}\right].
 \label{eq:exactmfull-mixture}
\end{align}
The positive factor \(a_\varepsilon\) has uniformly bounded first two
logarithmic time derivatives for all sufficiently small \(\varepsilon\).
\end{lemma}

\begin{proof}
Independence factors the contact event from every longitudinal slab
expectation.  By Eq.~\eqref{eq:exactmfull-stationary-variance}, convolution
of \(\phi_{j,\varepsilon}\) with the law of \(Z_t\) gives
\begin{equation}
 \mathbb E[\phi_{j,\varepsilon}(Z_t)]
 =\frac{1}{\sqrt{2\pi}\,\varepsilon S_*}
 \exp\!\left[-\frac{(\mu(t)-\mu(t_j))^2}
 {2\varepsilon^2S_*^2}\right].
 \label{eq:exactmfull-gaussian-convolution}
\end{equation}
Since
\(\mu(t)-\mu(t_j)=\varsigma\ell_0(x(t)-c_j)\),
Eqs.~\eqref{eq:exactmfull-common-width} and
\eqref{eq:exactmfull-gaussian-convolution} give
Eqs.~\eqref{eq:exactmfull-factorization}--
\eqref{eq:exactmfull-mixture}.  The last assertion follows from
lemma~\ref{lem:exactmfull-contact-tail}, because the remaining factors in
\(a_\varepsilon\) are independent of time.
\end{proof}

\subsection{The pure common-variance mixture}
\label{subsec:exactmfull-pure-mixture}

Set \(J=x(I)=[\alpha,\beta]\).  The targets are interior, so
\begin{equation}
 \alpha<c_1<\cdots<c_m<\beta.
 \label{eq:exactmfull-centre-order}
\end{equation}
For \(m\ge2\), let
\begin{equation}
 \Delta_j=c_{j+1}-c_j,
 \qquad \Delta_*:=\min_{1\le j<m}\Delta_j>0.
 \label{eq:exactmfull-gaps}
\end{equation}
No empty minimum is used when \(m=1\).

Define
\begin{equation}
 q_j(x)=w_j\e^{-(x-c_j)^2/(2\sigma^2)},
 \qquad
 \pi_j(x)=\frac{q_j(x)}{H_{\sigma,w}(x)},
 \qquad
 \bar c(x)=\sum_{j=1}^m\pi_j(x)c_j.
 \label{eq:exactmfull-posterior}
\end{equation}
The logarithmic slope \(L_{\sigma,w}=\partial_x\log H_{\sigma,w}\) obeys
the exact identities
\begin{align}
 L_{\sigma,w}(x)
 &=\frac{\bar c(x)-x}{\sigma^2},
 \label{eq:exactmfull-log-slope}\\
 L_{\sigma,w}'(x)
 &=\frac{\operatorname{Var}_{\pi(x)}(c)}{\sigma^4}
   -\frac{1}{\sigma^2}.
 \label{eq:exactmfull-log-slope-derivative}
\end{align}
Indeed, differentiation of \(q_j\) gives the first identity, while
\(\pi_j'=\pi_j(c_j-\bar c)/\sigma^2\) gives the second.

The mode cap for univariate Gaussian mixtures is classical
\cite{carreiraPerpinanWilliams2003modes,amendolaEngstromHaase2020modes}.
We retain the direct multiplicity-counting proof because the later Doi
transfer also needs explicit uniform root locations, curvature scales, and
full-complement margins.

\begin{lemma}[Global extended-Chebyshev zero bound]
\label{lem:exactmfull-zero-bound}
For distinct real centres and positive weights,
\(H_{\sigma,w}'\) has at most \(2m-1\) real zeros counted with
multiplicity.  Its real zero set is finite.
\end{lemma}

\begin{proof}
Multiplication by a nowhere-zero function preserves both zeros and their
multiplicities, and direct differentiation gives
\begin{equation}
 \sigma^2\e^{x^2/(2\sigma^2)}H_{\sigma,w}'(x)
 =\sum_{j=1}^m w_j(c_j-x)
   \e^{-c_j^2/(2\sigma^2)}\e^{c_jx/\sigma^2}.
 \label{eq:exactmfull-exponential-polynomial}
\end{equation}
We prove the slightly more general claim that a nonzero function
\begin{equation}
 P_m(x)=\sum_{j=1}^m(a_j+b_jx)\e^{\lambda_jx},
 \qquad \lambda_1<\cdots<\lambda_m,
 \label{eq:exactmfull-general-exp-polynomial}
\end{equation}
with no identically zero summand has at most \(2m-1\) real zeros counted
with multiplicity.

First, the zero set is finite.  After multiplication by
\(\e^{-\lambda_1x}\), the first affine term dominates as \(x\to-\infty\);
the last affine-exponential term dominates as \(x\to+\infty\).  Each
dominating affine factor has a fixed sign sufficiently far in its tail.
Thus all zeros lie in a compact interval.  A nonzero real-analytic function
has only finitely many zeros there.

Now induct on \(m\).  The \(m=1\) case is the affine zero bound.  For
\(m>1\), put \(Q=\e^{-\lambda_1x}P_m\).  Two derivatives annihilate the first
affine term.  For every \(j\ge2\), the second derivative of
\((a_j+b_jx)\e^{(\lambda_j-\lambda_1)x}\) is a nonzero affine polynomial
times the same exponential, and the remaining positive exponents are
distinct.  If a \(C^2\) function has \(N\ge2\) real zeros counted with
multiplicity, generalized Rolle counting gives at least \(N-2\) zeros of
its second derivative counted with multiplicity: differentiation reduces
the multiplicity at each root, and ordinary Rolle supplies a root between
each pair of distinct roots.  Hence
\begin{equation}
 N(Q)-2\le N(Q'')\le2(m-1)-1,
 \label{eq:exactmfull-rolle-count}
\end{equation}
so \(N(P_m)\le2m-1\); the cases \(N<2\) are immediate.  In
Eq.~\eqref{eq:exactmfull-exponential-polynomial}, every affine coefficient
has nonzero slope
\(-w_j\e^{-c_j^2/(2\sigma^2)}\), so the claim applies.
\end{proof}

\begin{lemma}[Uniform adjacent isolation]
\label{lem:exactmfull-adjacent-isolation}
For the fixed finite centre set and the compact weight family, there are
\(C_{\rm iso},q_{\rm iso}>0\) such that, for every \(j<m\),
\begin{equation}
 \sup_{x\in[c_j,c_{j+1}]}
 \frac{\sum_{k\notin\{j,j+1\}}q_k(x)}
      {q_j(x)+q_{j+1}(x)}
 \le C_{\rm iso}\e^{-q_{\rm iso}/\sigma^2}
 \label{eq:exactmfull-adjacent-isolation}
\end{equation}
for all sufficiently small \(\sigma\), uniformly in \(w\).  On
\([\alpha,c_1]\), component \(1\) has the analogous one-component bound, and
on \([c_m,\beta]\), component \(m\) does.
\end{lemma}

\begin{proof}
For \(x\in[c_j,c_{j+1}]\) and \(k<j\),
\begin{equation}
 (x-c_k)^2-(x-c_j)^2
 =(c_j-c_k)(2x-c_j-c_k)\ge(c_j-c_k)^2.
 \label{eq:exactmfull-left-square-comparison}
\end{equation}
For \(k>j+1\), comparison with \(c_{j+1}\) gives the corresponding positive
lower bound.  Since \(w_k/w_j\le1/w_*\), summation over the finite centre set
proves Eq.~\eqref{eq:exactmfull-adjacent-isolation}.  The same
difference-of-squares comparison proves the two tail statements.  The
posterior mean and variance of the full mixture therefore differ from those
of the isolated adjacent pair by \(O(\e^{-q/\sigma^2})\), uniformly in the
gap, the allocation, and sufficiently small \(\sigma\).
\end{proof}

\begin{theorem}[Exact topology of the pure mixture]
\label{thm:exactmfull-pure-topology}
Uniformly for \(w\in\mathcal W_{w_*}\), every sufficiently small
\(\sigma>0\) gives exactly \(m\) simple local maxima and \(m-1\) simple local
minima of \(H_{\sigma,w}\) on \(J\), ordered alternately, with nonzero
endpoint derivatives.  The maximum near \(c_j\), denoted \(p_j(\sigma,w)\),
satisfies
\begin{equation}
 p_j=c_j+O(\e^{-q/\sigma^2}),
 \qquad L_{\sigma,w}'(p_j)
 =-\sigma^{-2}+o(\sigma^{-2}).
 \label{eq:exactmfull-peak-location}
\end{equation}
For \(m\ge2\), define the weighted equality point
\begin{equation}
 s_j(\sigma,w)
 =\frac{c_j+c_{j+1}}2
  +\frac{\sigma^2}{\Delta_j}\log\frac{w_j}{w_{j+1}}.
 \label{eq:exactmfull-weighted-crossover}
\end{equation}
The minimum in the \(j\)-th gap, denoted \(r_j(\sigma,w)\), satisfies
\begin{align}
 r_j&=s_j+O(\sigma^4),
 \label{eq:exactmfull-valley-location}\\
 L_{\sigma,w}'(r_j)
 &=\frac{\Delta_j^2}{4\sigma^4}+O(\sigma^{-2}).
 \label{eq:exactmfull-valley-curvature}
\end{align}
All remainders are uniform on the compact weight set.
\end{theorem}

\begin{proof}
Fix a constant \(A>1\).  On
\(|x-c_j|\le A\sigma^2\), one-component isolation gives
\begin{align}
 L_{\sigma,w}(x)
 &=\frac{c_j-x}{\sigma^2}
   +O(\sigma^{-2}\e^{-q/\sigma^2}),
 \label{eq:exactmfull-peak-slope-expansion}\\
 L_{\sigma,w}'(x)
 &=-\sigma^{-2}
   +O(\sigma^{-4}\e^{-q/\sigma^2}).
 \label{eq:exactmfull-peak-derivative-expansion}
\end{align}
At \(c_j-\sigma^2\) the slope is positive, at
\(c_j+\sigma^2\) it is negative, and it is strictly decreasing between
them.  Hence there is one simple maximum with the location and derivative in
Eq.~\eqref{eq:exactmfull-peak-location}.

For a gap, the isolated adjacent odds are exactly
\begin{equation}
 \frac{q_{j+1}(x)}{q_j(x)}
 =\exp\!\left[\frac{\Delta_j(x-s_j)}{\sigma^2}\right].
 \label{eq:exactmfull-adjacent-odds}
\end{equation}
Put \(h_j=(\log9)/\Delta_j\) and
\begin{equation}
 C_j=[s_j-h_j\sigma^2,s_j+h_j\sigma^2].
 \label{eq:exactmfull-crossover-interval}
\end{equation}
At the left and right endpoints the adjacent posterior masses are,
respectively, \((9/10,1/10)\) and \((1/10,9/10)\).  By
lemma~\ref{lem:exactmfull-adjacent-isolation}, after reducing a common
\(\sigma_0\), both adjacent posterior masses are bounded below throughout
\(C_j\), and
\begin{equation}
 \operatorname{Var}_{\pi(x)}(c)\ge\kappa_0\Delta_*^2,
 \qquad
 L_{\sigma,w}'(x)\ge\frac{\kappa_v}{\sigma^4}>0
 \quad(x\in C_j)
 \label{eq:exactmfull-crossover-monotonicity}
\end{equation}
for constants independent of \(j,w,\sigma\).  At the two endpoints of
\(C_j\), Eq.~\eqref{eq:exactmfull-log-slope} gives, for sufficiently small
\(\sigma\),
\begin{equation}
 L_{\sigma,w}(s_j-h_j\sigma^2)
 \le-\frac{\Delta_j}{5\sigma^2},
 \qquad
 L_{\sigma,w}(s_j+h_j\sigma^2)
 \ge \frac{\Delta_j}{5\sigma^2}.
 \label{eq:exactmfull-crossover-signs}
\end{equation}
Thus \(C_j\) contains exactly one simple minimum root.

At \(s_j\), the adjacent posterior masses are equal and the adjacent
posterior mean is \((c_j+c_{j+1})/2\).  Hence
\begin{equation}
 L_{\sigma,w}(s_j)
 =-\frac{1}{\Delta_j}\log\frac{w_j}{w_{j+1}}
 +O(\sigma^{-2}\e^{-q/\sigma^2})=O(1)
 \label{eq:exactmfull-slope-at-crossover}
\end{equation}
uniformly in \(w\).  The mean-value theorem and
Eq.~\eqref{eq:exactmfull-crossover-monotonicity} give
\(r_j-s_j=O(\sigma^4)\).  At \(r_j\), the two adjacent posterior masses are
\(1/2+O(\sigma^2)\); substitution in
Eq.~\eqref{eq:exactmfull-log-slope-derivative} proves
Eq.~\eqref{eq:exactmfull-valley-curvature}.

For small \(\sigma\), the constructed peak and gap neighbourhoods are
pairwise disjoint and contained in \((\alpha,\beta)\).  They contain
\(m+(m-1)=2m-1\) distinct simple stationary points.  The global bound in
lemma~\ref{lem:exactmfull-zero-bound} excludes every additional real
stationary point, including multiple ones.  Tail isolation and the fixed
separations \(c_1-\alpha>0\) and \(\beta-c_m>0\) give
\begin{equation}
 L_{\sigma,w}(\alpha)>0,
 \qquad L_{\sigma,w}(\beta)<0,
 \label{eq:exactmfull-endpoint-signs}
\end{equation}
so neither endpoint is stationary.  When \(m=1\), there are no gaps or
valley definitions: the exact identity
\(L_{\sigma,w}(x)=(c_1-x)/\sigma^2\) gives one maximum and both endpoint
signs directly.
\end{proof}

\subsection{Full posterior-sector certificate and a slow positive factor}
\label{subsec:exactmfull-slow-factor}

The zero bound above does not apply after multiplication by a nonconstant
positive factor.  We therefore certify the complete complement of small peak
and valley boxes.

\begin{lemma}[Posterior-sector certificate]
\label{lem:exactmfull-posterior-sectors}
Fix \(K>0\).  There are constants
\(A_{\rm p},A_{\rm v},C_{\rm p},C_{\rm v},
\kappa_{\rm p},\kappa_{\rm v}>0\) and \(\sigma_K>0\), depending only on the
fixed centre set, \(J\), \(w_*\), and \(K\), such that, for every
\(0<\sigma<\sigma_K\) and every admissible \(w\), the boxes
\begin{align}
 P_j&=[c_j-A_{\rm p}\sigma^2,c_j+A_{\rm p}\sigma^2],
 \label{eq:exactmfull-peak-boxes}\\
 V_j&=[r_j-A_{\rm v}\sigma^4,r_j+A_{\rm v}\sigma^4]
 \quad(j<m)
 \label{eq:exactmfull-valley-boxes}
\end{align}
are pairwise disjoint and contained in \((\alpha,\beta)\), and the following
hold:
\begin{align}
 |L_{\sigma,w}|&\le C_{\rm p},&
 L_{\sigma,w}'&\le-\kappa_{\rm p}\sigma^{-2}
 &&\text{on }P_j,
 \label{eq:exactmfull-peak-box-control}\\
 |L_{\sigma,w}|&\le C_{\rm v},&
 L_{\sigma,w}'&\ge \kappa_{\rm v}\sigma^{-4}
 &&\text{on }V_j.
 \label{eq:exactmfull-valley-box-control}
\end{align}
At the left and right box boundaries, respectively, the slope signs are
\begin{equation}
 \begin{array}{c|cc}
 &\text{left boundary}&\text{right boundary}\\ \hline
 P_j&L_{\sigma,w}\ge2K&L_{\sigma,w}\le-2K\\
 V_j&L_{\sigma,w}\le-2K&L_{\sigma,w}\ge2K.
 \end{array}
 \label{eq:exactmfull-box-boundary-signs}
\end{equation}
On the entire complement in \(J\) of the open boxes,
\begin{equation}
 |L_{\sigma,w}(x)|\ge2K,
 \label{eq:exactmfull-complement-margin}
\end{equation}
with the exhaustive alternating sign order
\begin{equation}
 +\;P_1\;-\;V_1\;+\;P_2\;-\;\cdots\;+\;P_m\;-.
 \label{eq:exactmfull-complement-sign-order}
\end{equation}
For \(m=1\), every valley clause is omitted.
\end{lemma}

\begin{proof}
Choose \(A_{\rm p}\ge4K\).  The uniform expansions
\eqref{eq:exactmfull-peak-slope-expansion}--
\eqref{eq:exactmfull-peak-derivative-expansion} give, after reducing
\(\sigma_K\), Eq.~\eqref{eq:exactmfull-peak-box-control} and the peak row of
Eq.~\eqref{eq:exactmfull-box-boundary-signs}.  The convex-hull property
\(\bar c(x)\in[c_1,c_m]\) gives the full outer-tail estimates without any
local expansion:
\begin{align}
 x\le c_1-A_{\rm p}\sigma^2&\Longrightarrow
 L_{\sigma,w}(x)\ge A_{\rm p},\nonumber\\
 x\ge c_m+A_{\rm p}\sigma^2&\Longrightarrow
 L_{\sigma,w}(x)\le-A_{\rm p}.
 \label{eq:exactmfull-outer-tail-sector}
\end{align}

Fix a gap and abbreviate \(\Delta=\Delta_j\), \(s=s_j\).  The right
adjacent posterior mass is
\begin{equation}
 p(x)=\frac{q_{j+1}(x)}{q_j(x)+q_{j+1}(x)}
 =\frac{1}{1+\exp[-\Delta(x-s)/\sigma^2]}.
 \label{eq:exactmfull-logistic-posterior}
\end{equation}
On the crossover interval \(C_j\) in
Eq.~\eqref{eq:exactmfull-crossover-interval}, its edge odds are \(1/9\) and
\(9\), not exponentially small or large.  Both adjacent masses are therefore
bounded below throughout \(C_j\), and
Eq.~\eqref{eq:exactmfull-crossover-monotonicity} holds there.  Conversely,
the fixed finite centre diameter bounds the posterior variance, so
\(|L'|\le C\sigma^{-4}\) on \(C_j\).

The pure root satisfies \(L(r_j)=0\).  Choose \(A_{\rm v}\) large enough that
\(\kappa_vA_{\rm v}\ge2K\).  Since
\(A_{\rm v}\sigma^4=o(\sigma^2)\), the valley box lies in \(C_j\) for small
\(\sigma\).  Integrating the lower and upper bounds for \(L'\) from \(r_j\)
proves Eq.~\eqref{eq:exactmfull-valley-box-control}, the valley row of
Eq.~\eqref{eq:exactmfull-box-boundary-signs}, and the required negative and
positive bounds on \(C_j\setminus\operatorname{int}V_j\).

It remains to cover both outer sectors of every gap; this is the step that
prevents an unproved dominance claim at a crossover edge.  Consider first
\begin{equation}
 [c_j+A_{\rm p}\sigma^2,s_j-h_j\sigma^2]
 \label{eq:exactmfull-left-gap-sector}
\end{equation}
and split it at \(c_j+\Delta_j/4\), which lies in the stated interval for all
sufficiently small \(\sigma\), uniformly in \(w\).  On the portion nearer
\(c_j\), every other component is exponentially small relative to component
\(j\), so
\begin{equation}
 \bar c(x)-x\le-\tfrac12A_{\rm p}\sigma^2,
 \qquad L_{\sigma,w}(x)\le-A_{\rm p}/2\le-2K.
 \label{eq:exactmfull-left-near-peak-sector}
\end{equation}
On the remaining portion, monotonicity of
Eq.~\eqref{eq:exactmfull-logistic-posterior} and the left crossover edge give
\(p(x)\le1/10\), while \(x-c_j\ge\Delta_j/4\).  The adjacent-pair mean then
obeys
\begin{equation}
 c_j+\Delta_jp(x)-x
 \le\frac{\Delta_j}{10}-\frac{\Delta_j}{4}
 =-\frac{3\Delta_j}{20}.
 \label{eq:exactmfull-left-crossover-sector}
\end{equation}
The nonadjacent correction is exponentially small by
lemma~\ref{lem:exactmfull-adjacent-isolation}; hence this sector also has
\(L\le-2K\) for small \(\sigma\), including its \(1/9\)-odds edge.

Symmetrically, split
\([s_j+h_j\sigma^2,c_{j+1}-A_{\rm p}\sigma^2]\) at
\(c_{j+1}-\Delta_j/4\).  On the portion nearer the crossover,
\(p(x)\ge9/10\) and \(c_{j+1}-x\ge\Delta_j/4\), giving the positive analogue
of Eq.~\eqref{eq:exactmfull-left-crossover-sector}; on the portion nearer
\(c_{j+1}\), one-component isolation gives
\(L\ge A_{\rm p}/2\ge2K\).  Thus the right outer sector has \(L\ge2K\).

The fixed centre and endpoint separations make all boxes disjoint and
interior for one common sufficiently small \(\sigma_K\).  The outer tails,
all peak boundaries, both outer sectors of every gap, and every
crossover-minus-valley sector cover every point of the box complement.
This proves Eqs.~\eqref{eq:exactmfull-complement-margin}--
\eqref{eq:exactmfull-complement-sign-order}.  When \(m=1\), the same result
follows directly from \(L=(c_1-x)/\sigma^2\) and
Eq.~\eqref{eq:exactmfull-outer-tail-sector}.
\end{proof}

\begin{theorem}[Slow positive factors preserve the exact topology]
\label{thm:exactmfull-slow-factor}
Let \(x\in C^2(I)\) be strictly increasing with
\(\inf_Ix'=v_0>0\).  Let \(a_\sigma\in C^2(I)\) be positive and assume
\begin{equation}
 M_0:=\sup_{\sigma<\sigma_0}
 \|\partial_t\log a_\sigma\|_{L^\infty(I)}<\infty,
 \qquad
 M_1:=\sup_{\sigma<\sigma_0}
 \|\partial_t^2\log a_\sigma\|_{L^\infty(I)}<\infty.
 \label{eq:exactmfull-slow-factor-hypothesis}
\end{equation}
Then, uniformly for \(w\in\mathcal W_{w_*}\), all sufficiently small
\(\sigma>0\) give exactly \(m\) nondegenerate maxima and \(m-1\)
nondegenerate minima of
\begin{equation}
 F_{\sigma,w}(t)=a_\sigma(t)H_{\sigma,w}(x(t))
 \label{eq:exactmfull-slow-product}
\end{equation}
on \(I\), ordered alternately, with no stationary endpoint.  Relative to the
pure roots in the \(x\) coordinate,
\begin{align}
 x(t_{j,\sigma}^{\max})-p_j(\sigma,w)&=O(\sigma^2),
 \label{eq:exactmfull-slow-peak-shift}\\
 x(t_{j,\sigma}^{\min})-r_j(\sigma,w)&=O(\sigma^4),
 \label{eq:exactmfull-slow-valley-shift}
\end{align}
uniformly in \(w\).  Thus
\begin{equation}
 x(t_{j,\sigma}^{\min})=s_j(\sigma,w)+O(\sigma^4).
 \label{eq:exactmfull-slow-valley-location}
\end{equation}
\end{theorem}

\begin{proof}
Put
\begin{equation}
 b_\sigma=\partial_t\log a_\sigma,
 \qquad
 D_{\sigma,w}(t)=\partial_t\log F_{\sigma,w}(t)
 =b_\sigma(t)+x'(t)L_{\sigma,w}(x(t)).
 \label{eq:exactmfull-stationarity-function}
\end{equation}
Since \(F_{\sigma,w}>0\), its stationary points are precisely the zeros of
\(D_{\sigma,w}\), and
\begin{equation}
 D_{\sigma,w}'
 =b_\sigma'+x''L_{\sigma,w}+(x')^2L_{\sigma,w}'.
 \label{eq:exactmfull-stationarity-derivative}
\end{equation}
Choose \(K>(M_0+1)/v_0\) and apply
lemma~\ref{lem:exactmfull-posterior-sectors}.  On the full complement of the
open boxes, the term \(x'L\) strictly dominates \(b_\sigma\), so \(D\) has
the alternating signs in Eq.~\eqref{eq:exactmfull-complement-sign-order} and
has no complement zero.

On a peak box, Eqs.~\eqref{eq:exactmfull-peak-box-control} and
\eqref{eq:exactmfull-stationarity-derivative} give
\begin{equation}
 D_{\sigma,w}'
 \le M_1+\|x''\|_\infty C_{\rm p}
 -v_0^2\kappa_{\rm p}\sigma^{-2}<0
 \label{eq:exactmfull-peak-root-monotonicity}
\end{equation}
for small \(\sigma\).  Its boundary signs give exactly one zero.  At that
zero \(F''=FD'<0\), hence the root is a nondegenerate maximum.  Similarly,
on a valley box,
\begin{equation}
 D_{\sigma,w}'
 \ge-M_1-\|x''\|_\infty C_{\rm v}
 +v_0^2\kappa_{\rm v}\sigma^{-4}>0,
 \label{eq:exactmfull-valley-root-monotonicity}
\end{equation}
so its unique zero is a nondegenerate minimum.  The certified tail signs
exclude endpoint stationary points.

At a slow-factor root, Eq.~\eqref{eq:exactmfull-stationarity-function} gives
\(|L|\le M_0/v_0\).  The pure peak root has \(L(p_j)=0\), and the uniform
negative \(L'=\Theta(\sigma^{-2})\) on its box implies the peak displacement
in Eq.~\eqref{eq:exactmfull-slow-peak-shift}.  The uniform positive
\(L'=\Theta(\sigma^{-4})\) on a valley box gives
Eq.~\eqref{eq:exactmfull-slow-valley-shift}.  Combining the latter with
Eq.~\eqref{eq:exactmfull-valley-location} proves
Eq.~\eqref{eq:exactmfull-slow-valley-location}.  Every constant is uniform
over the compact allocation family.
\end{proof}

\subsection{Fixed-\texorpdfstring{\(\varepsilon\)}{epsilon} Doi transfer and
theorem boundary}
\label{subsec:exactmfull-doi-transfer}

For \(B>0\), let
\begin{equation}
 f_{B,\varepsilon,w}(t)
 =\mathbb E_0\!\left[
 K_{B,w,\varepsilon}(Z_t,R_t)
 \exp\!\left(-\int_0^tK_{B,w,\varepsilon}(Z_s,R_s)\,\dd s\right)
 \right]
 \label{eq:exactmfull-reaction-density}
\end{equation}
be the Doi reaction-time density and set
\begin{equation}
 \mathcal F_{B,\varepsilon,w}=f_{B,\varepsilon,w}/B.
 \label{eq:exactmfull-rescaled-density}
\end{equation}
This is budget-rescaling, not normalization to unit time integral.

\begin{theorem}[Exact \(m\) Doi modes for fixed finite \((d,m)\)]
\label{thm:exactmfull-doi}
Fix all data in Eqs.~\eqref{eq:exactmfull-midpoint-sde}--
\eqref{eq:exactmfull-killing}, including the stationary midpoint variance,
mutual independence, the strict weighted-space variance inequalities, target
times strictly inside \(I\), consistently oriented physical centres,
normalized slabs of physical width \(\varepsilon\rho\), the whole-window
contact margin, and the compact allocation family.  Then there is
\(\varepsilon_0>0\), depending on these fixed data, such that for every
\(0<\varepsilon<\varepsilon_0\) and every
\(w\in\mathcal W_{w_*}\), the continuum free-exposure clock
\(G_{\varepsilon,w}\) has exactly \(m\) nondegenerate maxima and \(m-1\)
nondegenerate minima on \(I\), ordered alternately, and no stationary
endpoint.

For each fixed such \(\varepsilon\), there is
\(B_0(\varepsilon)>0\), uniform over \(w\in\mathcal W_{w_*}\), such that
\(\mathcal F_{B,\varepsilon,w}\) has the same complete stationary signature
for every \(0<B<B_0(\varepsilon)\).  Multiplication by \(B>0\) changes no
stationary point, so the statement also holds for the physical density
\(f_{B,\varepsilon,w}\).  This is theorem~\ref{thm:main-modality} of
section~\ref{sec:exact-m-spine}, with its hypotheses written out in full.
\end{theorem}

\begin{proof}
In lemma~\ref{lem:exactmfull-factorization}, take
\(a_\sigma=a_\varepsilon\), with
\(\sigma=\varepsilon S_*/\ell_0\).  Lemma
\ref{lem:exactmfull-contact-tail} supplies
Eq.~\eqref{eq:exactmfull-slow-factor-hypothesis}; hence
theorem~\ref{thm:exactmfull-slow-factor} proves the exact stationary
signature of \(G_{\varepsilon,w}\), uniformly over the compact allocation
family.

Fix one \(0<\varepsilon<\varepsilon_0\).  The ordered simple roots depend
continuously on \(w\) by the implicit-function theorem.  Their finitely many
graphs over the compact set \(\mathcal W_{w_*}\) are compact and disjoint.
Continuity of the second time derivative gives pairwise disjoint root tubes
on which \(G''\) has a uniform nonzero sign and magnitude.  After shrinking
the tubes if necessary, \(G'\) has fixed opposite signs at the two boundaries
of each tube.  On the compact complement in
\(I\times\mathcal W_{w_*}\), \(|G'|\) has a positive minimum; the two endpoint
slopes have positive uniform absolute minima as well.

For this fixed \(\varepsilon\), \(V_{w,\varepsilon}\) is bounded uniformly in
\(w\), and Eqs.~\eqref{eq:exactmfull-midpoint-initial}--
\eqref{eq:exactmfull-relative-perp-initial} put \(q_0\) in the required
weighted state space.  Because \(w_j\ge w_*>0\), the compact allocation set
lies strictly inside the simplex and admits the affine control chart and a
complex tube required by theorem~\ref{thm:mixed-jet}.  (If the set is the
singleton \(w_j=1/m\), the same conclusion is the corresponding pointwise
case.)  Applying that theorem with time orders \(0,1,2\) and no control
derivative gives
\begin{equation}
 \sup_{w\in\mathcal W_{w_*}}
 \|\mathcal F_{B,\varepsilon,w}-G_{\varepsilon,w}\|_{C^2(I)}
 \longrightarrow0
 \qquad(B\downarrow0).
 \label{eq:exactmfull-fixed-epsilon-c2-transfer}
\end{equation}
Choose \(B_0(\varepsilon)\) so that the \(C^1\) error is smaller than all
boundary, complement, and endpoint margins and the second-derivative error
is smaller than the root-tube curvature margin.  In each tube,
\(\mathcal F'\) has opposite boundary signs and \(\mathcal F''\) keeps the
sign of \(G''\), so it has exactly one root of the same type.  The complement
and endpoints contain none.  This proves the complete finite-window
signature.
\end{proof}

\begin{remark}[Quantifiers, saturation, and scope]
\label{rem:exactmfull-scope}
The quantifier order is
\begin{equation}
 \text{fix }(d,m,\text{all geometric and probabilistic data})
 \ \Longrightarrow\
 0<\varepsilon<\varepsilon_0
 \ \Longrightarrow\
 0<B<B_0(\varepsilon).
 \label{eq:exactmfull-quantifier-order}
\end{equation}
There is no asserted lower bound for \(B_0(\varepsilon)\) uniform as
\(\varepsilon\downarrow0\), and no interchange of the two limits.  The proof
permits \(B_0(\varepsilon)\) to shrink rapidly as \(\varepsilon\)
decreases---potentially exponentially in \(1/\varepsilon^2\)---and no rate is
claimed; this is why no common useful budget is asserted across the family.
The construction is not uniform in \(d\) or \(m\), does not provide one
geometry with arbitrarily many modes, does not count stationary points
outside \(I\), and gives no process-level event-mass floor.  It treats
normalized longitudinal Gaussian slabs, not arbitrary localized catalyst
patches.

Moreover, lemma~\ref{lem:exactmfull-contact-tail} shows that the relative
contact factor is asymptotically saturated on the theorem window:
\begin{equation}
 c_{d,\varepsilon}=1
 +O(\varepsilon^{-N}\e^{-q/\varepsilon^2})
 \quad\text{in }C^2(I).
 \label{eq:exactmfull-saturated-contact}
\end{equation}
Thus the exact Doi embedding is rigorous, but the modal basis is created by
the ordered narrow slabs sampled by the monotone constant-variance midpoint;
the relative encounter factor is a slow perturbation in this asymptotic
family.  Nontrivial contact, a useful common positive budget, deterministic
finite-parameter complement certificates, survival, and event-basin masses
remain separate continuum-validation requirements rather than consequences
of theorem~\ref{thm:exactmfull-doi}.
\end{remark}

\nocite{grebenkov2020multiple}
\bibliographystyle{iopart-num}
\bibliography{references}

\end{document}
